\section{Benchmarks in the judgmental-heuristics history}
\label{sec:annex-p19-heuristics-program-history}

These proofs establish the mathematical comparisons stated in
\cref{sec:p19-heuristics-program-history}. They do not establish that an
empirical participant used Bayes' rule, independent sampling or normalized
support, nor that any reviewed manipulation uniquely identifies a process.

\begin{annexitem}{annex:p19.hph.odds}{thm:p19-hph-odds}{Prior odds and likelihood ratios}
\begin{proof}
The hypotheses form a partition, and all denominators specified in
\cref{def:p19-hph-odds-domain} are positive. The multiplication identity
for conditional probability gives
\[
 P(H_i\cap D)=P(D\mid H_i)P(H_i),\qquad
 P(H_i\mid D)=\frac{P(D\mid H_i)P(H_i)}{P(D)}
 \quad(i=1,2).
\]
In particular \(P(H_2\mid D)>0\). Dividing the two equations cancels
\(P(D)\) and yields
\[
 \frac{P(H_1\mid D)}{P(H_2\mid D)}
 =\frac{P(D\mid H_1)}{P(D\mid H_2)}
  \frac{P(H_1)}{P(H_2)}=\ell o_0.
\]
Since \(H_1,H_2\) partition the space, their conditional probabilities
sum to one. Put \(x=P(H_1\mid D)\). The preceding equality says
\(x/(1-x)=\ell o_0\), with \(1-x>0\). Multiplying and rearranging
gives \(x=\ell o_0(1-x)\), hence
\(x=\ell o_0/(1+\ell o_0)\). This also covers the allowed case
\(P(D\mid H_1)=0\), where both \(\ell\) and \(x\) are zero.

For the first stipulated example, \(o_0=1/9\) and \(\ell=9\),
so posterior odds equal one and \(x=1/(1+1)=1/2\).
For the second, \(\ell=1/9\), so posterior odds equal \(1/81\)
and \(x=(1/81)/(1+1/81)=1/82\). The arithmetic changes no
sampling assumption and supplies no empirical prevalence estimate.
\end{proof}
\end{annexitem}

\begin{annexitem}{annex:p19.hph.patterns}{thm:p19-hph-patterns}{The source's sampling and choice comparisons}
\begin{proof}
Under independent births with the same girl probability \(p\), the
probability of a specified string is the product of its eight coordinate
probabilities. Four coordinates contribute \(p\), and the other four
contribute \(1-p\), giving \(p^4(1-p)^4\). Order does not change
that product. Since \(p\in(0,1)\), it is positive.

The positions of the four girls specify a four-element subset of eight
positions. There are
\[
 \binom84=\frac{8\cdot7\cdot6\cdot5}{4\cdot3\cdot2\cdot1}=70
\]
such subsets and hence 70 mutually exclusive strings. At \(p=1/2\),
each product is \((1/2)^8=1/256\). Summing over those disjoint
strings gives \(70/256\). More generally, the balanced-count event
has probability \(70p^4(1-p)^4>0\). Any specified balanced string
is contained in it, so its probability conditional on that event is
\(p^4(1-p)^4/(70p^4(1-p)^4)=1/70\). A specified alternating string
is one of the 70 strings, not their union. This proves every birth-order
claim and makes the conditioning denominator explicit.

For the independent choice/outcome comparison, let \(R\) denote the
red outcome and \(C\) the choice of red, with \(P(R)=q\) and
\(P(C)=a\). Success is the disjoint union
\((C\cap R)\cup(C^c\cap R^c)\). Independence gives
\[
 P(\text{success})=aq+(1-a)(1-q)
                    =(1-q)+a(2q-1).
\]
At \(q=.7\), this equals \(.3+.4a\), which is increasing on
\(a\in[0,1]\) and is maximized at \(a=1\) with value .70.
Matching \(a=.7\) gives \(.49+.09=.58\). Linearity of expectation
then gives the corresponding expected number of successes over any
fixed number of such trials; it does not require that every realized
sequence attains its expectation. Changing the objective or making
the chooser's action informative about the trial would change the
premises and would require a different calculation.

Finally, under a uniform one-bean draw from each supplied urn, all beans
in that urn are equally likely. Red probability is therefore its red
count divided by its total count: \(1/10=.10\), \(9/100=.09\),
and \(5/100=.05\). With the same positive reward \(r\) for red,
zero otherwise, and no differing cost, expected rewards are
\(.10r,.09r,.05r\), respectively. Their strict ordering proves
the stated benchmark. This does not infer anyone's subjective beliefs,
utility or process from an observed bowl choice.
\end{proof}
\end{annexitem}

\begin{annexitem}{annex:p19.hph.support}{thm:p19-hph-support}{Normalized support and changes in one score}
\begin{proof}
Put \(x=s(A)\), \(y=s(B)\). Nonnegativity and \(x+y>0\)
give \(0\leq x/(x+y)\leq1\). Reversing the compared descriptions
gives \(y/(x+y)\), so their sum is \((x+y)/(x+y)=1\).
This normalization is local to the two scores and implies no additivity
for other descriptions.

If \(x'>x\geq0\) and \(y>0\), both sums are positive and
\[
 \frac{x'}{x'+y}-\frac{x}{x+y}
   =\frac{y(x'-x)}{(x'+y)(x+y)}>0.
\]
If instead \(y'>y\geq0\) and \(x>0\), then
\[
 \frac{x}{x+y'}-\frac{x}{x+y}
   =\frac{-x(y'-y)}{(x+y')(x+y)}<0.
\]
These establish the strict monotonicity claims. When \(y=0\) and
\(x>0\), increasing \(x\) leaves its ratio at one. When \(x=0\)
and \(y>0\), increasing \(y\) leaves the ratio at zero. The
zero--zero pair remains excluded rather than being assigned a value.

The constructed packed example gives \(2/(2+2)=1/2\); changing
only the first score to four gives \(4/(4+2)=2/3\). No property of
the fraction guarantees that unpacking produces the score four, or
even that it increases a score. The empirical premise about the effect
of descriptions must be established separately. The displayed arithmetic
is therefore a conditional model comparison, not proof of a psychological
unpacking effect or extensional subadditivity of a probability law.
\end{proof}
\end{annexitem}
